\chapter{데이터 준비와 시점 학습}

\section{상태 정리: \texttt{data.py}}
\textbf{하려는 것.} 원본의 \ft{loan_status}는 ``Fully Paid'', ``Charged Off'', ``Late (31-120 days)'' 등 문자열이다. 이것을 상환(paid), 부도(default), 연체(late), 진행 중(current) 네 가지로 정리하고, 숫자로 쓸 열(금리, 만기, 발행월)을 만든다.
\begin{lstlisting}
def classify(status):
    s = status.astype(str)
    return pd.Series(np.select(
        [s.str.contains('Fully Paid'), s.str.contains('Charged Off') | s.eq('Default'),
         s.str.startswith('Late') | s.eq('In Grace Period')],
        ['paid', 'default', 'late'], 'current'), index=status.index)

def parse(df):
    d = df.copy()
    d['ym'] = pd.to_datetime(d['issue_d'], format='%b-%Y')
    d['issue_year'] = d['ym'].dt.year
    d['r'] = pd.to_numeric(d['int_rate'].astype('string').str.strip().str.rstrip('%')) / 100
    d['term_m'] = d['term'].astype('string').str.extract(r'(\d+)')[0].astype(int)
    lp = pd.to_datetime(d['last_pymnt_d'], format='%b-%Y', errors='coerce')
    d['tau'] = months_between(d['ym'], lp).fillna(1).clip(lower=1).astype(int)
    d['outcome'] = classify(d['loan_status'])
    d['resolved'] = d['outcome'].isin(['paid', 'default'])
    d['bad'] = (d['outcome'] == 'default').astype('int8')
    return d
\end{lstlisting}
\begin{itemize}[leftmargin=*]
\item \textbf{3--6번 줄} \ft{np.select([조건1, 조건2, 조건3], [값1, 값2, 값3], 기본값)}: 엑셀의 중첩 IF. 위에서부터 처음 참인 조건의 값을 쓰고, 아무것도 아니면 \ft{'current'}다. \ft{.str.contains}는 ``그 글자를 포함하나'', \ft{.str.startswith}는 ``그 글자로 시작하나''다.
\item \textbf{9번 줄} \ft{df.copy()}: 원본을 건드리지 않으려고 복사본에서 작업한다.
\item \textbf{10번 줄}: ``Dec-2015'' 같은 글자를 날짜로 바꾼다. \texttt{\%b}는 월 약자, \texttt{\%Y}는 네 자리 연도다. 이 \ft{ym}(발행월)이 이후 모든 ``발행월마다''의 기준이다.
\item \textbf{12번 줄}: ``13.56\%'' \ensuremath{\to} 공백 제거 \ensuremath{\to} \% 제거 \ensuremath{\to} 숫자 13.56 \ensuremath{\to} 100으로 나눠 0.1356.
\item \textbf{13번 줄}: `` 36 months''에서 숫자만 뽑는다. \texttt{r'(\textbackslash d+)'}는 ``연속된 숫자''라는 정규식이다.
\item \textbf{14--15번 줄} \ft{tau}: 발행부터 마지막 납부까지 개월 수. 현금흐름 복원(7장)과 시점 학습(아래)에 쓴다. \ft{errors='coerce'}는 ``변환이 안 되면 빈 값으로''이고, \ft{fillna(1)}은 빈 값을 1로 채운다.
\item \textbf{17--18번 줄}: 결과가 확정된(상환 또는 부도) 대출만 \ft{resolved}=참. 부도면 \ft{bad}=1이 모델의 정답이다.
\end{itemize}

\paragraph{누수 목록 \ft{LEAK}.} \ft{total_pymnt}(받은 총액), \ft{recoveries}(회수금), \ft{last_pymnt_d}처럼 \textbf{승인한 뒤에야 알 수 있는} 열들이다. \ft{feature_columns()}가 이것을 모델 입력에서 뺀다. 이 열들은 채점(IRR)에는 쓰지만 고르기(모델)에는 절대 쓰지 않는다.
\begin{warnbox}[누수]
\ft{total_pymnt}를 모델에 넣으면 AUC가 0.99 가까이 나온다. 이미 결과를 보고 맞히는 것이라 실전에서는 쓸 수 없다. 새 열을 넣을 때마다 ``승인하는 날 이 값을 알 수 있나?''를 묻는다.
\end{warnbox}

\section{시점 학습: \texttt{asof.py}}
\textbf{하려는 것.} 2016년 1월 1일에 승인 결정을 한다면, 그날까지 \textbf{결과가 확정된} 대출로만 배워야 한다. 2015년 9월에 나간 대출이 2016년 3월에 부도났다는 사실은 2016년 1월에는 모른다.
\begin{lstlisting}
CHARGE_OFF_LAG = 5

def months_to(d, cutoff):
    c = pd.Timestamp(cutoff)
    return ((c.year - d['ym'].dt.year) * 12 + c.month - d['ym'].dt.month).to_numpy()

def known_at(d, cutoff, lag=CHARGE_OFF_LAG):
    age = months_to(d, cutoff)
    settle = d['tau'].to_numpy() + np.where(d['bad'].to_numpy() == 1, lag, 0)
    return (age > 0) & d['resolved'].to_numpy() & (settle <= age)

def training_set(d, cutoff, lag=CHARGE_OFF_LAG):
    return d[(d['ym'] < pd.Timestamp(cutoff)).to_numpy() & known_at(d, cutoff, lag)]
\end{lstlisting}
\begin{itemize}[leftmargin=*]
\item \textbf{1번 줄}: 부도는 마지막 납부 뒤 약 150일(5개월) 연체 후에 상각(Charged Off)되므로, 5개월 늦게 확정된다고 본다.
\item \textbf{5번 줄} \ft{age}: 발행월부터 기준일까지 개월 수. (연도 차이 \ensuremath{\times} 12) + 월 차이.
\item \textbf{9번 줄} \ft{settle}: 결과가 확정된 개월. 상환이면 \ft{tau}, 부도면 \ft{tau} + 5.
\item \textbf{10번 줄}: 세 조건을 모두 만족해야 학습에 쓸 수 있다. 기준일 전에 발행됐고, 결과가 확정됐고, 확정 시점이 기준일 이전이다.
\end{itemize}
\begin{workout}[따라 해 보기: \ft{snippets.py asof}]
기준일 2016-01-01, 대출 4건:
\begin{lstlisting}[style=out]
          ym  tau  bad  resolved  age  settle  known
0 2014-01-01   12    0      True   24      12   True
1 2014-06-01    6    1      True   19      11   True
2 2015-03-01    3    0      True   10       3   True
3 2015-09-01    2    1      True    4       7  False
\end{lstlisting}
\Step{1} 3번 대출: 2015년 9월 발행, 2개월 납부 뒤 부도. 확정은 $2+5=7$개월째인데 기준일까지는 4개월뿐이다.
\Ans{3번은 2016년 1월에는 ``아직 모르는'' 대출이라 학습에서 빠진다. 이 규칙이 없으면 미래를 보고 배우는 셈이다.}
\end{workout}
\Trap{기준일 이전에 \emph{발행된} 대출을 모두 쓰는 것. 그러면 빨리 끝난 대출(대부분 조기상환)만 남거나, 아직 모를 부도가 섞인다. 둘 다 결과를 왜곡한다(보고서 3장의 종결 편향).}
